\section{Data and Leakage Audit}\label{sec:data}

\textbf{Elliptic++.} The dataset~\cite{elmougy2023} contains 203{,}769 Bitcoin
transactions in 49 time steps of about two weeks (January 2016 to October 2017),
234{,}355 fund-flow edges between transactions, 822{,}942 addresses, 477{,}117
address$\rightarrow$transaction (input) edges and 837{,}124
transaction$\rightarrow$address (output) edges. Of the transactions, 4{,}545
are labelled illicit, 42{,}019 licit and 157{,}205 unknown. Labels are derived
from the entity controlling the transaction's input addresses~\cite{weber2019,safar2026}.
We verified that none of the 234{,}355 transaction edges crosses a time step:
each step is a separate graph.

\textbf{Temporal protocol.} We follow the chronological split of
Weber~et~al.~\cite{weber2019}: steps 1--29 are used to fit models, steps 30--34
form a validation window used only for early stopping and for choosing the
decision threshold, and steps 35--49 form the untouched test set (16{,}670
labelled transactions, 1{,}083 illicit). Every model family uses exactly the
same windows. Standardisation statistics are computed on training-period
nodes only.

\textbf{Feature audit.} Using the feature semantics recovered
in~\cite{safar2026}, we partition the 182 transaction features of Elliptic++
by \emph{when} they become observable (Table~\ref{tab:features}). Address
statistics (lifetime, balance, transaction counts) were computed once at block
575{,}059, about a year and a half after the last test transaction, and aggregates
over \emph{successor} transactions depend on later spending. Both violate the
rule that a feature must be observable before the prediction
time~\cite{kaufman2012,kapoor2023}. We therefore define two feature sets:
\textsc{All}, as used in the literature, and \textsc{Clean}, which keeps only
transaction-intrinsic features and aggregates over predecessor transactions.
The original features are z-scored over the whole dataset; because this is an
affine map, re-standardising on the training window removes its effect.

\begin{table}[t]
\caption{Elliptic++ transaction features grouped by observability.}
\label{tab:features}
\centering
\footnotesize
\setlength{\tabcolsep}{3pt}
\resizebox{\columnwidth}{!}{%
\begin{tabular}{@{}llrc@{}}
\toprule
Group & Columns (Elliptic++ names) & \# & \textsc{Clean} \\
\midrule
Tx-intrinsic & Local 1--21 & 21 & \checkmark \\
Address statistics at block 575{,}059 & Local 22--93 & 72 & -- \\
Predecessor aggregates & Aggregate 1--36 & 36 & \checkmark \\
Successor aggregates & Aggregate 37--72 & 36 & -- \\
Elliptic++ tx additions & in-degree, BTC, fee, size, \ldots & 16 & \checkmark \\
Successor count & out\_txs\_degree & 1 & -- \\
\midrule
\multicolumn{2}{@{}l}{Total (\textsc{All} / \textsc{Clean})} & 182 & 73 \\
\bottomrule
\end{tabular}}
\end{table}

The 965 transactions that Elliptic++ could not de-anonymise lack the 16
additions; we impute training-period medians and add a missing-value flag.
We do not use the released Elliptic++ address features, which are lifetime
summaries with the same look-ahead problem~\cite{liang2026}.
